% Propietario conservado: /Users/ruben/Documents/New project/output/REVISION_ARGUMENTAL_APERTURAS_CIERRES_20260915/III/spanish_source/manuscrito/sections/03b_vacancias_prefactor.tex
% Ampliación propia de III; no modifica el núcleo común copiado de I.
% Residencia propuesta: antes de Retorno de hoja, incidencia areal--volumétrica.
% Procedencia y controles: VACANCIAS_PREFACTOR_PROCEDENCIA.md, en esta carpeta.
\section{Polarisation discrète et mode tritique des projections arithmétiques}
\label{iii:iii:sec:vacancias-prefactor}

Le raffinement et la comparaison entre les feuillets d’APP fournissent deux
observations complémentaires. La première suit l’écart accumulé
entre le nombre de transitions et la capacité acquise ; la seconde conserve
le mode du sélecteur tritique dans l’appariement des fibres additives et
multiplicatives. La polarisation discrète et le préfacteur électronique
s’obtiennent à partir de ces opérations respectives. Leur articulation avec le vide
requiert de maintenir tant l’information temporelle que les espaces sur
lesquels agissent les opérateurs.

\subsection{Discrépance du raffinement et convention des extrémités}

Les cardinaux du bloc ternaire et de sa représentation décimale fixent
\[
 \rho=\log_{1000}729=\log_{10}9,\qquad
 \delta=1-\rho=\log_{10}(10/9).
\]
On a $0<\delta<1$. En outre, $\delta$ est irrationnel : une égalité $\delta=p/q$, avec $q>0$, impliquerait $10^{q-p}=9^q$, incompatible avec la décomposition en facteurs premiers. Ce décalage intervient déjà dans le codage du raffinement. Sa lecture centrée est explicitée ici, en conservant l'origine de phase et la convention des extrémités.

\begin{definition}[Codages inférieur et supérieur]
Pour $\theta\in\mathbb R$ et $n\in\mathbb Z$, on définit
\begin{align}
 z_n^{\downarrow}(\theta)
 &=\lfloor(n+1)\delta+\theta\rfloor-\lfloor n\delta+\theta\rfloor,
 \label{iii:iii:vp:lower}\\
 z_n^{\uparrow}(\theta)
 &=\lceil(n+1)\delta+\theta\rceil-\lceil n\delta+\theta\rceil.
 \label{iii:iii:vp:upper}
\end{align}
Les deux prennent leurs valeurs dans $\{0,1\}$. La flèche distingue exclusivement la convention des extrémités des intervalles ; l'orientation tritique conserve sa définition précédente.
\end{definition}

Le codage inférieur coïncide avec celui utilisé dans le développement
du raffinement. Le codage supérieur permet de retrouver la formulation de la
polarisation avec la fonction plafond. Les deux expressions sont reliées par
un terme de frontière exactement déterminé.

\begin{proposition}[Changement de convention et terme de frontière]
Soit $b_n(\theta)=\mathbf1_{\mathbb Z}(n\delta+\theta)$. Alors
\begin{equation}
 z_n^{\uparrow}(\theta)-z_n^{\downarrow}(\theta)
 =b_n(\theta)-b_{n+1}(\theta).
 \label{iii:iii:vp:endpoint}
\end{equation}
En particulier, pour $N\geq0$,
\[
 \sum_{n=0}^{N-1}(z_n^{\uparrow}-z_n^{\downarrow})
 =b_0-b_N.
\]
\end{proposition}
\begin{proof}
Pour tout réel $x$, $\lceil x\rceil=\lfloor x\rfloor+1-
\mathbf1_{\mathbb Z}(x)$. Appliquer cette identité aux deux extrémités
de chaque accroissement donne \eqref{iii:iii:vp:endpoint} ; leur somme est télescopique.
L’irrationalité de $\delta$ implique que, $\theta$ étant fixée, au plus
un indice entier satisfait $n\delta+\theta\in\mathbb Z$.
\end{proof}

Avec $\theta=0$, le premier symbole supérieur est $1$ et l'inférieur est $0$ ; pour $n\geq1$, ils coïncident. Le terme initial appartient au registre et est conservé lors du changement de convention. Les densités et les bornes de discrépance sont transportées avec ce terme, sans modifier la mise à jour du TPK ni recommencer l'histoire à chaque retour de phase.

\subsection{Polarisation bornée et bilan temporel exact}

\begin{definition}[Polarisation discrète centrée]
La polarisation du codage supérieur, relativement à la densité
$\delta$, est
\begin{equation}
 P_N^{\uparrow}(\theta)
 =\sum_{n=0}^{N-1}\bigl(z_n^{\uparrow}(\theta)-\delta\bigr),
 \qquad P_0^{\uparrow}=0.
 \label{iii:iii:vp:polarization}
\end{equation}
Son accroissement temporel est
\begin{equation}
 J_N^{\uparrow}(\theta)
 =P_{N+1}^{\uparrow}(\theta)-P_N^{\uparrow}(\theta)
 =z_N^{\uparrow}(\theta)-\delta.
 \label{iii:iii:vp:current}
\end{equation}
\end{definition}

Le centrage sépare la croissance uniforme de densité $\delta$ de
l’information de phase. Ainsi, $P_N^{\uparrow}$ mesure un déséquilibre accumulé,
non le nombre total de vacances. L’accroissement $J_N^{\uparrow}$ conserve
l’impulsion locale et le fond par rapport auquel elle est évaluée.

\begin{theorem}[Borne uniforme et continuité discrète temporelle]
Pour $N\geq0$ et toute phase $\theta$,
\begin{equation}
 P_N^{\uparrow}(\theta)
 =\lceil N\delta+\theta\rceil-\lceil\theta\rceil-N\delta,
 \qquad |P_N^{\uparrow}(\theta)|<1.
 \label{iii:iii:vp:bound}
\end{equation}
Pour tous entiers $0\leq a<b$,
\begin{equation}
 \sum_{n=a}^{b-1}J_n^{\uparrow}(\theta)
 =P_b^{\uparrow}(\theta)-P_a^{\uparrow}(\theta),
 \qquad
 \left|\sum_{n=a}^{b-1}J_n^{\uparrow}(\theta)\right|<1.
 \label{iii:iii:vp:continuity}
\end{equation}
La densité de la suite des vacances est $\delta$.
\end{theorem}
\begin{proof}
La somme des accroissements de plafond de \eqref{iii:iii:vp:upper} produit la première identité. En écrivant $r(x)=\lceil x\rceil-x\in[0,1)$, on obtient $P_N^{\uparrow}=r(N\delta+\theta)-r(\theta)$ ; sa valeur absolue est inférieure à un. Sur l'intervalle $[a,b)$, on obtient de même $r(b\delta+\theta)-r(a\delta+\theta)$, ce qui prouve la seconde borne, plus précise que la somme de deux bornes individuelles. Enfin, $N^{-1}\sum_{n=0}^{N-1}z_n^{\uparrow}
=\delta+N^{-1}P_N^{\uparrow}\to\delta$.
\end{proof}

Le codage inférieur admet les mêmes démonstrations en remplaçant
$r(x)$ par $\lfloor x\rfloor-x\in(-1,0]$. En particulier,
$P_N^{\uparrow}-P_N^{\downarrow}=b_0-b_N$. L’écart borné
utilisé précédemment dans les lectures de Dirichlet et la polarisation
présente conservent donc une relation explicite d’origine et d’extrémités.
L’apériodicité n’introduit pas non plus de dérive séculaire de ce déséquilibre :
l’histoire peut se prolonger indéfiniment tandis que l’écart
centré reste uniformément borné.

\subsection{Réalisation dans les droites de charge et de temps}

Pour réaliser dimensionnellement le bilan, on spécifie une droite de charge
$\mathcal L_Q$ de section non nulle $u_Q$ et une suite croissante
d’instants $t_N$ dans une droite temporelle orientée. Avec
$\Delta t_N=t_{N+1}-t_N>0$, la charge centrée et le courant moyen de
chaque intervalle sont définis par
\begin{equation}
 Q_N=P_N^{\uparrow}u_Q,\qquad
 I_N=\frac{Q_{N+1}-Q_N}{\Delta t_N}
     =\frac{u_Q}{\Delta t_N}\bigl(z_N^{\uparrow}-\delta\bigr).
 \label{iii:iii:vp:dimensional}
\end{equation}
Ainsi, $Q_N\in\mathcal L_Q$ et $I_N\in\mathcal L_Q\otimes\mathcal L_T^{-1}$. La section $u_Q$ et les intervalles temporels sont des données de cette réalisation ; la construction arithmétique précédente a déjà déterminé $P_N^{\uparrow}$ et $J_N^{\uparrow}$.

\begin{proposition}[Conservation du bilan sous réalisation dimensionnelle]
Pour $a<b$,
\begin{equation}
 Q_b-Q_a=\sum_{n=a}^{b-1}\Delta t_n I_n,
 \qquad
 \left|\frac{Q_b-Q_a}{u_Q}\right|<1.
 \label{iii:iii:vp:dimensional-balance}
\end{equation}
Si tous les intervalles ont longueur $t_0$, les deux courants possibles
sont $-\delta u_Q/t_0$ et $(1-\delta)u_Q/t_0$.
\end{proposition}
\begin{proof}
Multiplier chaque identité de \eqref{iii:iii:vp:dimensional} par
$\Delta t_n$ et sommer fait télescoper les charges. La borne est
\eqref{iii:iii:vp:continuity} ; les deux valeurs locales s’obtiennent à partir de
$z_n^{\uparrow}\in\{0,1\}$.
\end{proof}

La continuité démontrée est temporelle et se rapporte à la charge centrée. Une réalisation spatiale incorpore en outre la localisation des charges, l'incidence entre cellules et l'identification des flux de frontière. Dans une description constitutive par canaux, le couplage est spécifié au moyen d'une application d'excitation $\iota:\mathcal L_Q\to\mathcal H_{\rm ch}\otimes\mathcal L_Q$ et d'une application d'observation $\ell:\mathcal H_{\rm ch}\otimes\mathcal L_Q\to\mathcal L_Q$, où $\mathcal H_{\rm ch}$ est une réalisation réelle de l'espace des canaux. Si les deux sont linéaires et indépendantes de l'indice temporel, la réponse d'un opérateur $\mathcal V$ fixe est transportée par
\[
 Q_N^{\rm resp}=\ell(\mathcal V\otimes I_{\mathcal L_Q})\iota(Q_N),
 \qquad
 I_N^{\rm resp}=\frac{Q_{N+1}^{\rm resp}-Q_N^{\rm resp}}{\Delta t_N}.
\]
La linéarité conserve le bilan télescopique. Cette composition
explicite les données d’une identification constitutive : la séquence
scalaire détermine l’impulsion, tandis que $\iota$, $\mathcal V$ et $\ell$
déterminent son excitation, son transport et son observation par canaux. Les
identités \eqref{iii:iii:vp:polarization}--\eqref{iii:iii:vp:dimensional-balance}
établissent la première construction indépendamment de ce choix de
réalisation. L’opérateur du vide est construit dans la
section~\ref{iii:iii:sec:constitutiva} sur ses canaux et ses projecteurs propres.

\subsection{Appariement des fibres additives et multiplicatives}

La seconde observation provient de la réalisation tridimensionnelle d’APP,
$\mathcal A_3=D_9^3$, avec une mesure uniforme et
$D_9=\{1,\ldots,9\}$. On conserve les observables
\[
 \Sigma(x)=\rho_9(x_1+x_2+x_3),\qquad
 \Pi(x)=\rho_9(x_1x_2x_3),
\]
où $\rho_9(n)=1+[(n-1)\bmod9]$ pour $n>0$. Dans
$\ell^2(\mathcal A_3)$, soient $P_{\Sigma,3}$ et $P_{\Pi,3}$ les
espérances conditionnelles sur leurs fibres. Par exemple,
\[
 (P_{\Sigma,3}f)(x)=
 \frac1{|\Sigma^{-1}(\Sigma(x))|}
 \sum_{y:\,\Sigma(y)=\Sigma(x)}f(y).
\]
Ce sont des projections orthogonales définies sur le même support. La matrice
d’incidences $N_{ab}=|\Sigma^{-1}(a)\cap\Pi^{-1}(b)|$ résulte de
l’énumération des $729$ triplets, en additionnant et en multipliant leurs coordonnées avant
d’appliquer $\rho_9$ :
\begin{equation}
 N=9\begin{pmatrix}
 0&1&1&0&1&1&0&1&4\\
 1&0&1&1&0&1&1&0&4\\
 1&0&2&0&1&2&0&0&3\\
 0&1&1&0&1&1&0&1&4\\
 1&0&1&1&0&1&1&0&4\\
 0&0&2&1&0&2&0&1&3\\
 0&1&1&0&1&1&0&1&4\\
 1&0&1&1&0&1&1&0&4\\
 0&1&2&0&0&2&1&0&3
 \end{pmatrix}.
 \label{iii:iii:vp:joint-matrix}
\end{equation}
Les lignes ont pour somme $81$ et les sommes des colonnes sont
\[
 (m_1,\ldots,m_9)=(36,36,108,36,36,108,36,36,297).
\]
L’appariement des indicatrices normalisées des fibres
est donc $C_{ab}=N_{ab}/\sqrt{81m_b}$. La normalisation recueille les
multiplicités produites par APP, non un choix de valeurs singulières.

\begin{proposition}[Mode singulier du sélecteur tritique]
Dans la carte canonique, le vecteur du sélecteur est $m_{\rm ph}=(1,-1,0)^3$. On a
\begin{equation}
 Cm_{\rm ph}=C^{\mathsf T}m_{\rm ph}=-\tfrac12m_{\rm ph}.
 \label{iii:iii:vp:trit-mode}
\end{equation}
Ainsi, la valeur singulière correspondante est $s=1/2$.
\end{proposition}
\begin{proof}
Dans les six colonnes où $m_{\rm ph}$ est non nul, $m_b=36$ et
$C_{ab}=N_{ab}/54$. Multiplier les lignes de
\eqref{iii:iii:vp:joint-matrix} par $(1,-1,0)^3$ produit
$-m_{\rm ph}/2$. Pour le produit transposé, les sommes signées
des colonnes $3,6,9$ sont nulles ; les six autres donnent les mêmes
valeurs $-m_{{\rm ph},b}/2$. Cela prouve les deux identités et, par
composition, $CC^{\mathsf T}m_{\rm ph}=m_{\rm ph}/4$.
\end{proof}

\subsection{Norme du commutateur et préfacteur électronique}

\begin{theorem}[Norme exacte d’incompatibilité arithmétique]
Les deux projections satisfont
\begin{equation}
 \bigl\|[P_{\Sigma,3},P_{\Pi,3}]\bigr\|=\frac{\sqrt3}{4}.
 \label{iii:iii:vp:prefactor}
\end{equation}
Le maximum est atteint dans le plan principal correspondant au mode
tritique de \eqref{iii:iii:vp:trit-mode}.
\end{theorem}
\begin{proof}
La matrice rationnelle $M=CC^{\mathsf T}$ a pour entrées $M_{ac}=\sum_bN_{ab}N_{cb}/(81m_b)$. Son spectre complet se vérifie au moyen des sous-espaces indépendants suivants : les vecteurs $(1,\ldots,1)$, $(1,-1,0)^3$ et $(1,1,-2)^3$ ont pour valeurs propres $1$, $1/4$ et $7/132$, respectivement ; le sous-espace des vecteurs à support dans $\{3,6,9\}$ dont la somme est nulle a pour dimension $2$ et pour valeur propre $1/18$ ; les vecteurs de somme nulle à support dans $\{1,4,7\}$ ou dans $\{2,5,8\}$ forment un noyau de dimension $4$. La substitution dans la matrice vérifie ces actions et les dimensions ont pour somme $9$. Par conséquent,
\[
 \operatorname{spec}M=
 \{1,1/4,7/132,1/18,1/18,0,0,0,0\}.
\]
Pour deux projections orthogonales, une valeur singulière $s\in(0,1)$ de
l’appariement détermine un plan principal dans lequel leurs matrices sont
\[
 P=\begin{pmatrix}1&0\\0&0\end{pmatrix},\qquad
 Q=\begin{pmatrix}s^2&s\sqrt{1-s^2}\\
 s\sqrt{1-s^2}&1-s^2\end{pmatrix}.
\]
Cette forme s’obtient en prenant un vecteur unitaire $u$ de la première
image avec $PQu=s^2u$ et en le complétant par
$(Qu-s^2u)/(s\sqrt{1-s^2})$. Leur commutateur est
$s\sqrt{1-s^2}\bigl(\begin{smallmatrix}0&1\\-1&0\end{smallmatrix}\bigr)$.
Les secteurs de valeur singulière $0$ ou $1$ contribuent par zéro.
Par conséquent, la norme totale est le maximum de
$\sqrt{\lambda(1-\lambda)}$ sur le spectre de $M$. Ce maximum
est $\sqrt{3/16}$, atteint en $\lambda=1/4$.
\end{proof}

La norme est le préfacteur qui intervient dans la composition électronique.
Son obtention conserve le support tridimensionnel, les fibres et le mode
du sélecteur ; l’évaluation scalaire ne remplace pas cette structure. En
particulier, connaître $\sqrt3/4$ permet de connaître la grandeur maximale de
l’incompatibilité, tandis que les projecteurs conservent les directions
dans lesquelles elle agit.

\clearpage
\subsection{Défaut quadratique et réponse des secteurs \(90/120\)}

Le mode tritique détermine simultanément
\begin{equation}
 1-s^2=\frac34=\frac{90}{120},\qquad s=\frac12.
 \label{iii:iii:vp:defect-ratio}
\end{equation}
La première expression est le défaut quadratique du plan principal d'APP. La dernière conserve les longueurs des deux secteurs du calendrier. La réponse constitutive utilise ces secteurs sur un autre objet : un canal contractant $q\in(0,1)$. Son évaluation est
\[
 r(q)=\frac{1-q^{90}}{1-q^{120}}.
\]
La relation entre le quotient des secteurs et cette réponse peut être
précisée sans identifier leurs opérateurs.

\begin{proposition}[Limite du quotient de réponse]
Para $0<q<1$,
\begin{equation}
 \frac34<r(q)<1,\qquad
 \lim_{q\uparrow1}r(q)=\frac34=1-s^2.
 \label{iii:iii:vp:response-limit}
\end{equation}
\end{proposition}
\begin{proof}
Avec $u=q^{30}\in(0,1)$, la factorisation de $1-u^3$ et
$1-u^4$ donne
$r(q)=(1+u+u^2)/(1+u+u^2+u^3)$. La limite en $u=1$ est
$3/4$. De plus,
$4(1+u+u^2)-3(1+u+u^2+u^3)
=(1-u)(3u^2+2u+1)>0$, et le dénominateur dépasse le numérateur
de $u^3>0$. On obtient les deux inégalités.
\end{proof}

La coïncidence exacte concerne le défaut du mode singulier et la limite
de la réponse des secteurs. À l’intérieur contractif, la réponse
dépend de $q$ et est strictement supérieure à $3/4$. Le développement
constitutif conserve ses deux canaux $q_\pm$ et ses projecteurs
$P_\pm$ ; les projecteurs $P_{\Sigma,3}$ et $P_{\Pi,3}$ conservent,
pour leur part, les fibres arithmétiques qui ont déterminé le préfacteur.
Une équivalence opératorielle entre les deux réalisations exige une application
entre leurs espaces qui entrelace les actions correspondantes.
Les équations \eqref{iii:iii:vp:trit-mode}, \eqref{iii:iii:vp:prefactor} et
\eqref{iii:iii:vp:response-limit} précisent l’information partagée et
maintiennent visibles les données de chaque réalisation.
